% ============================================================================
% LANGENTWURF LE-05c: Das Verb gustar (gefallen/mögen)
% Spanisch Klasse 7 - Sequenz 5: Mi tiempo libre
% ============================================================================
% Tier: 1 (vollständig)
% Einzelstunde: 5c
% Fachfarbe: #c41e3a (Karminrot)
% ============================================================================

\documentclass[11pt,a4paper]{article}

\usepackage[utf8]{inputenc}
\usepackage[T1]{fontenc}
\usepackage[ngerman]{babel}
\usepackage{geometry}
\usepackage{fancyhdr}
\usepackage{lastpage}
\usepackage{xcolor}
\usepackage{tcolorbox}
\usepackage{enumitem}
\usepackage{tabularx}
\usepackage{longtable}
\usepackage{array}
\usepackage{booktabs}
\usepackage{amssymb}

% ============================================================================
% KONFIGURATION
% ============================================================================

\newcommand{\fach}{Spanisch}
\newcommand{\fachfarbe}{c41e3a}
\newcommand{\jahrgangsstufe}{7}
\newcommand{\schulart}{Gesamtschule}
\newcommand{\stundenthema}{Das Verb gustar (gefallen/mögen)}
\newcommand{\reihenthema}{Mi tiempo libre}
\newcommand{\stundennummer}{5c}
\newcommand{\reihenstunden}{4}
\newcommand{\gerniveau}{A1}
\newcommand{\lehrwerk}{Qué pasa 1}
\newcommand{\lehrwerkseite}{S. 73-74}
\newcommand{\lerngruppengroesse}{24}

% ============================================================================
% LAYOUT
% ============================================================================
\geometry{left=2.5cm, right=2.5cm, top=2.5cm, bottom=2.5cm}
\definecolor{fach}{HTML}{\fachfarbe}
\definecolor{gray}{HTML}{666666}
\definecolor{lightgray}{HTML}{f5f5f5}
\definecolor{successgreen}{HTML}{2a7a4b}
\definecolor{warningorange}{HTML}{b35c00}
\definecolor{basis}{HTML}{2a7a4b}
\definecolor{standard}{HTML}{2c5aa0}
\definecolor{erweiterung}{HTML}{7b1fa2}

\tcbuselibrary{skins,breakable}

\newtcolorbox{infobox}[1][]{
    colback=lightgray,
    colframe=fach,
    fonttitle=\bfseries,
    title=#1,
    breakable
}

\newtcolorbox{scorebox}{
    colback=white,
    colframe=fach,
    fonttitle=\bfseries,
    title=Didaktik-Score,
    breakable
}

\pagestyle{fancy}
\fancyhf{}
\fancyhead[L]{\small\textcolor{gray}{\fach{} -- Jg. \jahrgangsstufe{} -- \stundenthema}}
\fancyhead[R]{\small\textcolor{gray}{LE-\stundennummer{} | Tier 1}}
\fancyfoot[C]{\small\thepage{} / \pageref{LastPage}}
\renewcommand{\headrulewidth}{0.4pt}

% Differenzierungsmarker
\newcommand{\B}{\textcolor{basis}{\textbf{[B]}}}
\newcommand{\Std}{\textcolor{standard}{\textbf{[S]}}}
\newcommand{\E}{\textcolor{erweiterung}{\textbf{[E]}}}

% ============================================================================
% DOKUMENT
% ============================================================================
\begin{document}

% ============================================================================
% DECKBLATT
% ============================================================================
\begin{titlepage}
    \centering
    \vspace*{2cm}
    {\large\textcolor{gray}{\schulart}}\\[2cm]
    {\Huge\textcolor{fach}{\textbf{Langentwurf}}}\\[0.5cm]
    {\large LE-\stundennummer{} | Tier 1 (vollständig)}\\[2cm]
    {\LARGE\textbf{\stundenthema}}\\[0.5cm]
    {\large Sequenz: \reihenthema}\\[2cm]
    \begin{tabular}{rl}
        \textbf{Fach:} & \fach{} (\gerniveau) \\[0.3cm]
        \textbf{Klasse:} & \jahrgangsstufe \\[0.3cm]
        \textbf{Lehrwerk:} & \lehrwerk, \lehrwerkseite \\[0.3cm]
        \textbf{Stunde:} & \stundennummer{} \\[0.3cm]
    \end{tabular}
    \vfill
\end{titlepage}

\tableofcontents
\newpage

% ============================================================================
% 1. BEDINGUNGSANALYSE
% ============================================================================
\section{Bedingungsanalyse}

\subsection{Lerngruppe}
\begin{tabular}{ll}
    Kursgröße: & \lerngruppengroesse{} SuS \\
    Jahrgangsstufe: & \jahrgangsstufe \\
    Sprachniveau: & \gerniveau{} (GER) -- Anfänger \\
    Fremdsprache: & 2. Fremdsprache (nach Englisch) \\
\end{tabular}

\subsection{Vorwissen aus Sequenz 5a/5b}
\begin{itemize}
    \item Freizeitaktivitäten: jugar al fútbol, escuchar música, ver la tele, leer libros...
    \item Sportarten und Hobbys benennen
    \item Bestimmte Artikel: el, la, los, las
    \item Erste Verben: ser, estar, tener, hacer
\end{itemize}

\subsection{Differenzierung (intern)}
\begin{tabular}{|l|p{3.5cm}|p{3.5cm}|p{3.5cm}|}
\hline
& \B{} \textbf{Basis} & \Std{} \textbf{Standard} & \E{} \textbf{Erweiterung} \\
\hline
\textbf{Ziel} & BR: gusta/gustan erkennen & MR: Sätze bilden & GY: Alle Personen + Transfer \\
\hline
\textbf{gustar} & gusta vs. gustan unterscheiden & Mit me/te/le anwenden & Alle Pronomen + Negation \\
\hline
\textbf{Produktion} & Lückentext ausfüllen & Eigene Sätze schreiben & Dialoge führen \\
\hline
\end{tabular}

\vspace{0.5em}
\textbf{WICHTIG:} Diese Marker sind NUR für die Lehrkraft!

% ============================================================================
% 2. SACHANALYSE
% ============================================================================
\section{Sachanalyse}

\subsection{Sprachliche Strukturen}

\textbf{Das Verb \textit{gustar} -- besondere Konstruktion:}

Im Spanischen funktioniert \textit{gustar} anders als im Deutschen. Wörtlich: ``Mir gefällt...'' statt ``Ich mag...''.

\begin{center}
\begin{tabular}{|l|l|l|}
\hline
\textbf{Indirektes Objektpronomen} & \textbf{Deutsch} & \textbf{Beispiel} \\
\hline
me & mir & Me gusta el fútbol. \\
te & dir & ¿Te gusta la música? \\
le & ihm/ihr/Ihnen & Le gusta leer. \\
nos & uns & Nos gustan los deportes. \\
os & euch & ¿Os gusta bailar? \\
les & ihnen/Ihnen & Les gusta el cine. \\
\hline
\end{tabular}
\end{center}

\textbf{Regel für gusta vs. gustan:}
\begin{itemize}
    \item \textbf{gusta} + Infinitiv: Me \textbf{gusta} jugar al fútbol.
    \item \textbf{gusta} + Singular-Nomen: Me \textbf{gusta} el fútbol.
    \item \textbf{gustan} + Plural-Nomen: Me \textbf{gustan} los deportes.
\end{itemize}

\subsection{Kontrastive Betrachtung}

\begin{tabular}{|l|p{5cm}|p{5cm}|}
\hline
& \textbf{Deutsch} & \textbf{Spanisch} \\
\hline
Struktur & Ich mag + Objekt & Mir gefällt + Subjekt \\
\hline
Verb-Konjugation & Nach ``ich/du/er...'' & Nach dem, was gefällt (Sg./Pl.) \\
\hline
Beispiel & Ich mag Fußball. & Me gusta el fútbol. (= Mir gefällt der Fußball.) \\
\hline
\end{tabular}

\textbf{Typische Interferenz:} *``Yo gusto el fútbol'' (falsche Übertragung aus dem Deutschen)

\subsection{Kulturelle Aspekte}
\begin{itemize}
    \item Freizeitaktivitäten in spanischsprachigen Ländern
    \item Beliebte Sportarten: Fußball, Baloncesto, Tenis
    \item Musik: Reggaeton, Flamenco, Pop Latino
\end{itemize}

% ============================================================================
% 3. DIDAKTISCHE ANALYSE
% ============================================================================
\section{Didaktische Analyse}

\subsection{Stellung in der Sequenz}
\begin{tabular}{|c|l|l|}
\hline
\textbf{Block} & \textbf{Thema} & \textbf{Status} \\
\hline
5a & Freizeitaktivitäten (Wortschatz) & abgeschlossen \\
5b & Sportarten und Hobbys & abgeschlossen \\
\textbf{5c} & \textbf{Das Verb gustar} & \textbf{diese Stunde} \\
5d & Mi tiempo libre (Anwendung/Test) & folgt \\
\hline
\end{tabular}

\subsection{Kompetenzziele (Rahmenplan MV)}
\begin{itemize}
    \item \textbf{Kommunikativ:} Sprechen (über Vorlieben sprechen), Schreiben
    \item \textbf{Sprachlich:} Grammatik (gustar-Konstruktion), Wortschatz (Freizeit)
    \item \textbf{Interkulturell:} Freizeitgestaltung in hispanischen Kulturen
    \item \textbf{Methodisch:} Regelentdeckung, Partnerarbeit
\end{itemize}

\subsection{Legitimation}
\textbf{Rahmenplan Spanisch MV (Kl. 7-10):}
\begin{itemize}
    \item Kompetenzbereich: Kommunikative Kompetenz -- Sprechen
    \item Standard: ``Die SuS können über ihre Vorlieben und Interessen sprechen und danach fragen.''
\end{itemize}

% ============================================================================
% 4. STUNDENZIELE
% ============================================================================
\section{Stundenziele}

\subsection{Grobziel}
Die SuS erweitern ihre \textbf{kommunikative Kompetenz im Bereich Sprechen}, indem sie mit dem Verb \textit{gustar} über Vorlieben sprechen und die besondere Konstruktion korrekt anwenden.

\subsection{Feinziele (operationalisiert)}
\begin{tabular}{|c|p{7.5cm}|c|c|}
\hline
\textbf{Nr.} & \textbf{Die SuS können...} & \textbf{AFB} & \textbf{Diff.} \\
\hline
FZ1 & die Konstruktion von \textit{gustar} \textbf{erklären} und von der deutschen Struktur unterscheiden & I & alle \\
\hline
FZ2 & \textit{gusta} (Singular/Infinitiv) und \textit{gustan} (Plural) korrekt \textbf{unterscheiden} und einsetzen & II & alle \\
\hline
FZ3 & Sätze mit \textit{me/te/le gusta(n)} \textbf{bilden} und auf Fragen nach Vorlieben antworten & II & \Std{}/\E{} \\
\hline
FZ4 & die gustar-Konstruktion auf alle Personen \textbf{übertragen} und in einem Dialog \textbf{anwenden} & III & \E{} \\
\hline
\end{tabular}

\vspace{1em}
\textbf{AFB-Verteilung:} AFB I: 25\% | AFB II: 50\% | AFB III: 25\%

% ============================================================================
% 5. METHODISCHE ANALYSE
% ============================================================================
\section{Methodische Analyse}

\subsection{Einstieg}
\begin{itemize}
    \item \textbf{Methode:} Bildimpuls (Freizeitaktivitäten) + Lehreraussage: ``Me gusta el fútbol. ¿Y a ti?''
    \item \textbf{Begründung:} Aktivierung Vorwissen, Neugier wecken
    \item \textbf{Alternative:} Kurzes Video mit Jugendlichen über Hobbys
\end{itemize}

\subsection{Erarbeitung}
\begin{itemize}
    \item \textbf{Methode:} Entdeckendes Lernen -- gustar-Konstruktion aus Beispielen ableiten
    \item \textbf{Begründung:} Aktive Konstruktion statt passiver Rezeption
    \item \textbf{Scaffolding:} Kontrastierung Deutsch-Spanisch, visuelle Darstellung
\end{itemize}

\subsection{Übungsphasen}
\begin{itemize}
    \item \textbf{Methode:} Zuordnung (gusta/gustan) $\rightarrow$ Lückentext $\rightarrow$ Freie Produktion
    \item \textbf{Progression:} Rezeptiv $\rightarrow$ Reproduktiv $\rightarrow$ Produktiv
    \item \textbf{Differenzierung:} Hilfen für \B{}, Erweiterung für \E{}
\end{itemize}

\subsection{Medien}
\begin{itemize}
    \item Tafel/Whiteboard: gustar-Konstruktion visualisieren
    \item Lehrwerk: \lehrwerk, \lehrwerkseite
    \item Arbeitsblatt: AB-05c.pdf
    \item Lernpfad: LP-05c.html (Tablets, optional)
    \item Bildkarten: Freizeitaktivitäten
\end{itemize}

% ============================================================================
% 6. VERLAUFSPLANUNG
% ============================================================================
\section{Verlaufsplanung}

\textbf{Einzelstunde (45 Min.)}

\begin{longtable}{|p{1cm}|p{1.5cm}|p{4cm}|p{3cm}|p{1.5cm}|p{2cm}|}
\hline
\textbf{Zeit} & \textbf{Phase} & \textbf{Lehrerhandeln} & \textbf{Schülerhandeln} & \textbf{SF} & \textbf{Material} \\
\hline
\endhead

0--5 & Einstieg & Bilder zeigen: ``Me gusta el fútbol. ¿Qué te gusta a ti?'' & Vermutungen, erstes Sprechen & Plenum & Bilder \\
\hline

5--10 & Aktivier. & Beispiele sammeln: ``Me gusta..., Me gustan...'' -- Was fällt auf? & Muster erkennen & Plenum & Tafel \\
\hline

10--20 & Input & gustar-Konstruktion erklären: Pronomen (me, te, le...) + gusta/gustan & Tabelle ergänzen, Notizen & Plenum & Tafel, AB \\
\hline

20--28 & Übung 1 & EA: Zuordnung gusta vs. gustan (AB Aufg. 1-2) & AB bearbeiten & EA & AB-05c \\
\hline

28--38 & Übung 2 & PA: Dialoge ``¿Qué te gusta?'' -- ``Me gusta...'' & Dialoge üben & PA & Bildkarten \\
\hline

38--43 & Sicherung & 2-3 Paare präsentieren, Feedback & Präsentieren, zuhören & Plenum & -- \\
\hline

43--45 & Abschluss & Merkkasten sichern, HA: Vokabeln + LP-05c & Notieren & Plenum & -- \\
\hline
\end{longtable}

\textbf{Legende:} EA = Einzelarbeit, PA = Partnerarbeit

% ============================================================================
% 7. ERWARTETE SCHWIERIGKEITEN
% ============================================================================
\section{Erwartete Schwierigkeiten}

\begin{tabular}{|p{4.5cm}|p{8cm}|}
\hline
\textbf{Schwierigkeit} & \textbf{Reaktion} \\
\hline
Interferenz *``Yo gusto...'' & Kontrastierung: ``Mir gefällt'' vs. ``Ich mag'' -- Visualisierung an Tafel \\
\hline
Verwechslung gusta/gustan & Faustregel: Was gefällt? Singular = gusta, Plural = gustan \\
\hline
Pronomen-Verwechslung (me/te/le) & Tabelle als Hilfe, Pronomen farbig markieren \\
\hline
Vergessen des Artikels & Hinweis: ``Me gusta \textbf{el} fútbol'' -- Artikel gehört zum Nomen! \\
\hline
Zeitproblem & Puffer: Aufgabe 4 als HA aufgeben \\
\hline
\end{tabular}

% ============================================================================
% 8. HAUSAUFGABE
% ============================================================================
\section{Hausaufgabe}

\begin{itemize}
    \item AB-05c: Aufgabe 4 fertigstellen (freie Sätze)
    \item Vokabeln lernen: gustar-Konstruktion, Freizeitaktivitäten
    \item Optional: Lernpfad LP-05c.html bearbeiten
\end{itemize}

% ============================================================================
% 9. LÖSUNGEN
% ============================================================================
\section{Lösungen}

\subsection{AB-05c Lösungen}

\textbf{Merkkasten: gustar-Konstruktion}
\begin{itemize}
    \item me, te, le, nos, os, les
    \item gusta + Infinitiv/Singular
    \item gustan + Plural
\end{itemize}

\textbf{Aufgabe 1: gusta oder gustan?}
\begin{enumerate}[label=\alph*)]
    \item Me \textbf{gusta} el fútbol. (Singular)
    \item Me \textbf{gustan} los deportes. (Plural)
    \item Te \textbf{gusta} leer libros. (Infinitiv)
    \item Nos \textbf{gustan} las películas. (Plural)
    \item Le \textbf{gusta} la música. (Singular)
\end{enumerate}

\textbf{Aufgabe 2: Sätze vervollständigen}
\begin{enumerate}[label=\alph*)]
    \item ¿Te \textbf{gusta} jugar al fútbol?
    \item A María \textbf{le gusta} escuchar música.
    \item Nos \textbf{gustan} los videojuegos.
    \item ¿Os \textbf{gusta} bailar?
\end{enumerate}

\textbf{Aufgabe 3: Dialog}
\begin{itemize}
    \item A: ¿Qué te gusta hacer?
    \item B: Me gusta jugar al fútbol. ¿Y a ti?
    \item A: A mí me gusta escuchar música.
    \item B: ¿Te gustan los videojuegos?
    \item A: Sí, me gustan mucho.
\end{itemize}

\textbf{Aufgabe 4: Freie Produktion (Bewertungskriterien)}
\begin{itemize}
    \item Korrekte gustar-Konstruktion: 3P
    \item Richtige Wahl gusta/gustan: 2P
    \item Verwendung verschiedener Pronomen: 1P
\end{itemize}

% ============================================================================
% 10. ANLAGEN
% ============================================================================
\section{Anlagen}

\begin{enumerate}
    \item Arbeitsblatt (AB-05c.pdf)
    \item Musterlösung (ML-05c.pdf)
    \item Lehrerhinweise (LH-05c.pdf)
    \item Lernpfad (LP-05c.html)
    \item Bildkarten Freizeitaktivitäten (Kopiervorlage)
\end{enumerate}

% ============================================================================
% 11. DIDAKTIK-SCORE
% ============================================================================
\newpage
\section{Didaktik-Score}

\begin{scorebox}
\textbf{Bewertung der didaktischen Qualität nach 10 Dimensionen}\\
\textbf{Ziel-Score: $\geq$ 9.0 (Premium-Qualität)}

\vspace{1em}
\begin{tabular}{|c|l|c|c|p{4cm}|}
\hline
\textbf{\#} & \textbf{Dimension} & \textbf{Gew.} & \textbf{Score} & \textbf{Notiz} \\
\hline
1 & Differenzierung & 15\% & 9/10 & [B]/[S]/[E] qualitativ \\
\hline
2 & Taxonomie (AFB) & 10\% & 9/10 & 25/50/25 erfüllt \\
\hline
3 & Lernziel-Alignment & 10\% & 10/10 & RP MV explizit \\
\hline
4 & Aktivierung & 10\% & 9/10 & Entdeckend, Dialoge \\
\hline
5 & Scaffolding & 10\% & 9/10 & Kontrastierung DE-ES \\
\hline
6 & Feedback-Qualität & 10\% & 9/10 & LP: elaboriert \\
\hline
7 & Sprachliche Klarheit & 10\% & 10/10 & Kl. 7 angemessen \\
\hline
8 & Motivation & 10\% & 9/10 & Alltagsbezug Freizeit \\
\hline
9 & Inklusion (UDL) & 10\% & 9/10 & Mehrere Zugänge \\
\hline
10 & Praktikabilität & 5\% & 9/10 & 45 Min. realistisch \\
\hline
\multicolumn{3}{|r|}{\textbf{GESAMT:}} & \textbf{9.1/10} & \\
\hline
\end{tabular}

\vspace{1em}
$\checkmark$ \textcolor{successgreen}{\textbf{FREIGABE} (Score $\geq$ 9.0)}
\end{scorebox}

\end{document}
